\documentclass[11pt,a4paper]{article}
\usepackage[utf8]{inputenc}
\usepackage[T1]{fontenc}
\usepackage{lmodern}
\usepackage[french]{babel}
\usepackage{amsmath,amssymb,mathtools}
\usepackage{icomma}
\usepackage{xcolor}
\usepackage{colortbl}
\usepackage{tikz}
\usepackage[most]{tcolorbox}
\usepackage{enumitem}
\usepackage{array}
\usepackage{booktabs}
\usepackage{fancyhdr}
\usepackage[hidelinks]{hyperref}
\usetikzlibrary{arrows.meta,positioning,calc,patterns,backgrounds,shapes.geometric,decorations.pathreplacing,decorations.markings}
\usepackage[a4paper,margin=2.15cm,headheight=26pt,headsep=12pt]{geometry}
\usepackage{titlesec}

% ---------------- Palette ----------------
\definecolor{primary}{RGB}{0,151,169}
\definecolor{primarydark}{RGB}{0,88,101}
\definecolor{defcol}{RGB}{0,114,178}
\definecolor{thmcol}{RGB}{0,140,120}
\definecolor{formulacol}{RGB}{198,93,9}
\definecolor{excol}{RGB}{0,135,90}
\definecolor{attncol}{RGB}{200,40,55}
\definecolor{methcol}{RGB}{90,90,100}
\definecolor{remarkcol}{RGB}{120,94,200}
\definecolor{barcol}{RGB}{0,151,169}
\definecolor{barcold}{RGB}{0,110,124}
\definecolor{modecol}{RGB}{200,55,55}
\definecolor{meancol}{RGB}{0,90,200}
\definecolor{medcol}{RGB}{160,90,190}
\definecolor{quartcol}{RGB}{0,135,90}
\definecolor{ogivecol}{RGB}{176,74,0}

% ---------------- Boîtes ----------------
\tcbset{boxbase/.style={enhanced,breakable,boxrule=0.9pt,arc=2.5pt,
  left=3mm,right=3mm,top=2.3mm,bottom=2.3mm,
  fonttitle=\bfseries,coltitle=white,toptitle=0.6mm,bottomtitle=0.6mm}}

\newtcolorbox{definition}[1][]{boxbase,colback=defcol!6,colframe=defcol,
  title={\textsc{Définition}\ifx\relax#1\relax\relax\else\textmd{ : \itshape #1}\fi}}
\newtcolorbox{propriete}[1][]{boxbase,colback=thmcol!7,colframe=thmcol,
  title={\textsc{Propriété}\ifx\relax#1\relax\relax\else\textmd{ : \itshape #1}\fi}}
\newtcolorbox{formule}[1][]{boxbase,colback=formulacol!7,colframe=formulacol,
  title={\textsc{Formule}\ifx\relax#1\relax\relax\else\textmd{ : \itshape #1}\fi}}
\newtcolorbox{exemple}[1][]{boxbase,colback=excol!6,colframe=excol,
  title={\textsc{Exemple}\ifx\relax#1\relax\relax\else\textmd{ : \itshape #1}\fi}}
\newtcolorbox{methode}[1][]{boxbase,colback=methcol!8,colframe=methcol,sharp corners,
  title={\textsc{Méthode}\ifx\relax#1\relax\relax\else\textmd{ : \itshape #1}\fi}}
\newtcolorbox{remarque}[1][]{boxbase,colback=remarkcol!7,colframe=remarkcol,
  title={\textsc{Astuce}\ifx\relax#1\relax\relax\else\textmd{ : \itshape #1}\fi}}
\newtcolorbox{attention}[1][]{boxbase,colback=attncol!7,colframe=attncol,
  title={\textsc{Attention}\ifx\relax#1\relax\relax\else\textmd{ : \itshape #1}\fi}}

% Boîte "exercice" et "corrigé" (titre libre passé en argument)
\newtcolorbox{exobox}[1]{boxbase,colback=primary!4,colframe=primary,
  before skip=8pt,after skip=8pt,title={#1}}
\newtcolorbox{corbox}[1]{boxbase,colback=excol!5,colframe=excol!85!black,
  before skip=8pt,after skip=8pt,title={#1}}

% ---------------- En-tête ----------------
\newcommand{\docpart}[1]{\def\thedocpart{#1}}
\docpart{}
\pagestyle{fancy}
\fancyhf{}
\fancyhead[L]{\small\textcolor{primarydark}{\textbf{Fiche 8} — Statistiques descriptives}}
\fancyhead[R]{\small\textcolor{primarydark}{\textsc{\thedocpart}}}
\fancyfoot[C]{\small\thepage}
\renewcommand{\headrulewidth}{0.8pt}
\renewcommand{\headrule}{\hbox to\headwidth{\color{primary}\leaders\hrule height \headrulewidth\hfill}}

\titleformat{\section}{\large\bfseries\color{primarydark}}{\thesection}{0.6em}{}
\titleformat{\subsection}{\normalsize\bfseries\color{primary}}{\thesubsection}{0.6em}{}
\setlist{topsep=2pt,itemsep=1.5pt,parsep=0pt}

% Bandeau de titre du document
\newcommand{\bandeau}[2]{%
\begin{tcolorbox}[boxbase,colback=primary,colframe=primarydark,halign=center,
   before skip=2pt,after skip=12pt]
{\LARGE\bfseries\color{white} #1}\\[3pt]{\normalsize\color{white} #2}
\end{tcolorbox}}

\newcommand{\ecm}{\sigma_m}
\newcommand{\ect}{\sigma}
\newcommand{\med}{M\!e}
\docpart{Tutoriel — Thème 1}
\begin{document}
\bandeau{Thème 1 — Le tableau statistique}{Effectifs, fréquences, cumuls et centres de classe}

\noindent\textbf{But du thème.} Avant \emph{tout} calcul d'indicateur (moyenne, médiane, variance\ldots),
il faut une chose : un \textbf{tableau statistique propre et complet}. C'est la « matière première »
de toute la fiche. Ce thème entraîne la seule mécanique de construction de ce tableau, colonne par
colonne. Rien de difficile : de l'addition et une division. Mais tout le reste en dépend.

\section{Les cinq colonnes à savoir remplir}

\begin{formule}[les grandeurs de base]
Pour chaque ligne $i$ (une modalité si le caractère est discret, une classe $[a,b[$ s'il est continu) :
\[
\underbrace{x_i=\frac{a+b}{2}}_{\text{centre de classe}}\qquad
\underbrace{N=\sum_{i=1}^{k} n_i}_{\text{effectif total}}\qquad
\underbrace{f_i=\frac{n_i}{N}}_{\text{fréquence}}
\]
\begin{itemize}
  \item $n_i$ : \textbf{effectif} de la ligne = nombre d'individus concernés (un entier) ;
  \item $x_i$ : \textbf{centre de classe} = milieu de $[a,b[$ ; pour un caractère discret, $x_i$ est
        directement la valeur (pas de calcul) ;
  \item $N$ : \textbf{effectif total} = somme de toute la colonne $n_i$ ;
  \item $f_i$ : \textbf{fréquence} = proportion (entre $0$ et $1$) ; $f_i\times 100$ = pourcentage.
\end{itemize}
\end{formule}

\begin{exemple}[centre de classe]
Classe $[85,95[$ : $x=\dfrac{85+95}{2}=90$. \quad Classe $[0,4[$ : $x=\dfrac{0+4}{2}=2$.
Le centre est \emph{toujours} le milieu, même si les bornes sont « bizarres » :
$[12{,}5\,;17{,}5[\ \to\ x=\dfrac{12{,}5+17{,}5}{2}=15$.
\end{exemple}

\section{Les colonnes « cumulées »}

Cumuler, c'est \textbf{additionner de proche en proche}, du haut vers le bas.

\begin{definition}[effectif cumulé $n_c$ et fréquence cumulée $f_c$]
\[
n_c^{(i)} = n_1+n_2+\cdots+n_i \qquad\text{et}\qquad f_c^{(i)}=f_1+f_2+\cdots+f_i .
\]
$n_c^{(i)}$ = nombre d'individus dont la valeur est \emph{inférieure} à la borne supérieure de la
ligne $i$. La \textbf{dernière} valeur cumulée vaut toujours $n_c=N$ et $f_c=1{,}00$ ($=100\,\%$).
\end{definition}

\begin{methode}[construire le tableau en 5 gestes]
\begin{enumerate}
  \item Colonne \textbf{centres} $x_i=(a+b)/2$ (ou recopier les valeurs si discret).
  \item Colonne \textbf{effectifs} $n_i$, puis \textbf{additionner} pour obtenir $N$.
  \item Colonne \textbf{fréquences} $f_i=n_i/N$ (arrondir à $0{,}01$).
  \item Colonne \textbf{effectifs cumulés} $n_c$ : on descend en ajoutant à chaque fois $n_i$.
  \item Colonne \textbf{fréquences cumulées} $f_c$ : idem avec les $f_i$.
\end{enumerate}
\textbf{Deux contrôles obligatoires :} la somme de la colonne $n_i$ doit redonner $N$, et la somme
de la colonne $f_i$ doit valoir $1$ (dernier $f_c=1{,}00$).
\end{methode}

\begin{exemple}[une ligne se lit d'un coup d'œil]
\renewcommand{\arraystretch}{1.2}
\begin{center}
\rowcolors{2}{primary!5}{white}
\begin{tabular}{|c|c|c|c|c|c|}
\hline
\rowcolor{primary!18}
Classes & $x_i$ & $n_i$ & $f_i$ & $n_c$ & $f_c$ \\\hline
$[0,10[$   & 5  & 8  & 0{,}20 & 8  & 0{,}20 \\
$[10,20[$  & 15 & 14 & 0{,}35 & 22 & 0{,}55 \\
$[20,30[$  & 25 & 18 & 0{,}45 & 40 & 1{,}00 \\\hline
\multicolumn{2}{|r|}{Totaux} & $N=40$ & $1{,}00$ & & \\\hline
\end{tabular}
\end{center}
\emph{Lecture :} $22$ individus ($55\,\%$) ont une valeur inférieure à $20$. On vérifie
$8+14+18=40=N$ et $0{,}20+0{,}35+0{,}45=1{,}00$ : le tableau est cohérent.
\end{exemple}

\begin{attention}[pièges classiques]
\begin{itemize}
  \item Ne jamais confondre \textbf{effectif} $n_i$ (un compte, entier) et \textbf{fréquence} $f_i$
        (une proportion, entre $0$ et $1$).
  \item Une fréquence cumulée ne \emph{redescend jamais} : elle croît de $0$ jusqu'à $1$.
  \item Le centre de classe utilise \textbf{les deux bornes}, pas seulement la borne inférieure.
  \item Un caractère \textbf{discret} n'a pas de colonne « classes » : $x_i$ = la valeur elle-même.
\end{itemize}
\end{attention}

\begin{remarque}[retrouver un effectif à partir d'une fréquence]
La relation $f_i=\dfrac{n_i}{N}$ se lit dans les deux sens : si l'on connaît $f_i$ et $N$, alors
$n_i=f_i\times N$. Par exemple $f_i=0{,}15$ pour $N=200$ donne $n_i=0{,}15\times 200=30$ individus.
On s'en sert dès qu'un énoncé donne des pourcentages au lieu des effectifs.
\end{remarque}
\end{document}
